\section{引言}

深度卷积神经网络（CNN）自AlexNet在2012年ImageNet竞赛中取得突破性成果以来，一直是图像识别领域的主导技术。直觉上，网络的深度（层数）对于性能至关重要：更深的网络能够学习更丰富的层次化特征，从而提升表达能力。然而，随着层数不断增加，研究人员发现了一个令人困惑的现象：当网络深度增加到一定程度时，训练误差和测试误差不仅没有下降，反而同时上升。这种问题并非由过拟合引起（因为训练误差同样升高），而是属于优化上的困难。文献中将此称为\textbf{退化问题（Degradation Problem）}。

退化问题的存在严重限制了深层网络的实用价值。例如，一个56层的传统卷积网络在CIFAR-10数据集上的表现甚至不如一个20层的浅层网络，这违背了“越深越好”的直观假设。针对这一困境，He等人于2015年提出了\textbf{残差学习框架（Residual Learning）}，通过引入快捷连接（Shortcut Connection）来显式地构造恒等映射，使得网络可以轻松地学习到恒等映射附近的扰动。这一创新极大简化了超深网络的训练过程，成功训练了多达152层的网络（是当时VGG网络深度的8倍），并在ILSVRC 2015竞赛中斩获图像分类、检测、定位等多个项目的冠军。

本文旨在对ResNet论文的核心思想进行深度解析。我们将首先剖析退化问题的本质，然后详细介绍残差学习的数学原理与实现方式，接着展示网络架构的设计细节和丰富的实验验证，最后从梯度流动、集成学习等角度探讨残差网络为什么有效，并回顾其对后续深度学习发展的深远影响。通过本文的梳理，读者将能全面理解残差学习是如何使“深度成为可能”的。